% !TeX program = xelatex
\documentclass[11pt,a4paper]{article}
\usepackage[margin=23mm]{geometry}
\usepackage{fontspec}
\setmainfont{Latin Modern Roman}
\usepackage{amsmath,amssymb,mathtools}
\usepackage{xcolor,enumitem,booktabs,longtable,array,xurl,hyperref}
\definecolor{muted}{HTML}{53636C}
\hypersetup{colorlinks=true,urlcolor=blue,linkcolor=blue}
\setlength{\parindent}{0pt}
\setlength{\parskip}{5pt}
\newcommand{\R}{\mathbb R}
\newcommand{\E}{\mathbb E}
\newcommand{\N}{\mathbb N}
\newcommand{\ind}{\mathbf 1}
\newcommand{\Var}{\operatorname{Var}}
\newcommand{\Cov}{\operatorname{Cov}}
\newcommand{\supp}{\operatorname{supp}}
\DeclareMathOperator*{\argmin}{arg\,min}
\DeclareMathOperator*{\argmax}{arg\,max}
\newcommand{\entrymeta}[3]{{\sffamily\small\color{muted}\textbf{\texttt{#1}}\quad|\quad #2\par #3\par}}
\newcommand{\Problemref}[1]{\texttt{#1}}
\newcommand{\JELref}[2]{\texttt{#2} (JEL \texttt{#1})}
\begin{document}
\section*{Identifiable habit dynamics and competing state dependence}
\textbf{Problem ID:} \texttt{D15-229}\quad \textbf{JEL:} \texttt{D15}\par
% BEGIN ECONOMICS STATEMENT
\label{problem:OP-0588}
\label{beyond:DT-03}
\entrymeta{DT-03}{Behavioral decision theory}{Research agenda formalization; open-problem attribution unverified}
\subsection*{Definitions and assumptions}
Consider a consumer choosing an adapted consumption path subject to known intertemporal budgets. Its preference model is
\[
 \E\left[\sum_{t=0}^{T}\beta^tu(c_t,h_t;\theta)\right],\qquad
 h_{t+1}=G(h_t,c_t;\gamma),
\]
where $\beta\in(0,1)$ and $h_t$ is a habit state unobserved by the econometrician. Each model comparison must declare the consumption and state spaces with their measurable structures, parameter spaces, admissible utility and transition families, and finite horizon $T\in\{0,1,2,\ldots\}$; if several horizons are observed, their collection must be declared. The parameters $\theta,\gamma$ and initial habit are unknown. Specify whether the initial habit is a fixed unknown $h_0$ or has an unknown distribution, allowing a point mass to represent the former case. Declare whether $\beta$ is known or unknown, the income/price and shock processes, feasible budget sets, and the consumer's information filtration, including whether the consumer observes habit. Specify the econometrician's observations and the measurable choice rule, including any errors, randomization, tie-breaking or approximation rule. These primitives must determine a unique observable probability law; uniqueness of an optimizer is not assumed merely from the objective.

Write $m$ for this complete structural specification, including its nuisance primitives. For a declared common family $\mathcal Z$ of assigned experimental environments, let $P_m^z$ denote the induced law of the observed consumption path under environment $z$. This replaces the abbreviated indexing $P_{u,G,\theta,\gamma}(c_{0:T}\mid z)$: $P_m^z$ is an intervention-indexed law, not a conditional distribution at an arbitrary probability-zero event. Each environment must specify assignment, timing, anticipation and information, feasible consumption, price or income perturbations, and which structural mechanisms and initial-state laws are maintained across environments. A formulation using observational conditioning instead must specify an assignment law and compare conditional laws almost everywhere under that common law, unless particular versions are part of the experiment. Competing model families must be declared, for example time-separable utility with persistent unobserved types or a specified adjustment-cost family, together with their common experimental family and maintained nuisance restrictions.

Two structural models are observationally equivalent for the declared experimental family if $P_{m_1}^z=P_{m_2}^z$ for every $z\in\mathcal Z$. A permitted latent-state relabeling $\phi:H\to H'$ is a bijection for which $\phi$ and $\phi^{-1}$ are measurable and which preserves all additional structure required by the declared model class, including any smoothness restrictions. With $h'=\phi(h)$, the transformed primitives are
\[
 G^\phi(h',c;\gamma)=\phi\bigl(G(\phi^{-1}(h'),c;\gamma)\bigr),\qquad
 u^\phi(c,h';\theta)=u(c,\phi^{-1}(h');\theta).
\]
The initial state becomes $\phi(h_0)$, or its distribution becomes the image of the original initial-state distribution under $\phi$; all other state-dependent information and choice components must be transformed consistently. When both specifications are admissible, this change of coordinates leaves behavior unchanged. Identification of a latent habit technology therefore requires a declared normalization or a target defined modulo the permitted relabelings. A canonical representative must be specified rather than presumed to exist uniquely.
\subsection*{Formal question}
Find behavioral axioms, experimental interventions and substantive restrictions under which a canonical habit model is identified and observationally distinguishable from the specified heterogeneity and adjustment-cost alternatives. Characterize which impulse responses to exogenous consumption, price or income perturbations identify habit persistence and which remain confounded. Establish testable restrictions or prove observational equivalence within each explicitly defined model comparison.

This is a research programme over explicitly specified comparisons, rather than one theorem with a model class supplied in advance. Each comparison must state a nontrivial admissible class $\mathcal M$, a target $\tau(m)$ and its domain, and the horizon collection being observed. Point identification means that equality of all the declared observable laws for $m_1,m_2\in\mathcal M$ implies $\tau(m_1)=\tau(m_2)$; any local or generic version must state its neighborhood or exceptional-set convention. Specify whether the target is a parameter, utility, transition law, persistence functional or equivalence class, and whether functional identification is restricted to state-consumption pairs reachable under the declared experiments. Any canonical normalization must address the admitted observational invariances, not just state coordinates; for ordinary expected-utility maximization, positive affine utility transformations also preserve choices. A numerical persistence functional must be well defined under the admitted normalizations or invariant on the chosen equivalence classes. Identification obtained solely by fixing the target in advance is not the requested characterization.

For each rival comparison, characterize the intersection of the families' sets of observable experimental laws and state any exclusions used for distinguishability; membership in a shared observable law cannot itself distinguish the two families. Define each impulse response by its observed outcome functional, reference environment, perturbation magnitude or derivative, lag, population or conditioning event, and anticipation convention, with the required moments or derivatives justified. Here testable restrictions mean restrictions on population observable laws. A finite-sample testing formulation is an additional specialization requiring a sampling experiment, size criterion and power against a declared separated alternative class. The inherited common conventions do not supply undeclared local structural, causal or sampling assumptions.
\subsection*{Scope and status}
Habit is itself a form of state dependence; it cannot be distinguished from an unrestricted model allowed to reproduce the same history-dependent preferences. Restrictions on rival models are consequently essential. The attached catalogue attributes a habit research direction to an interview with Colin Camerer. The publicly accessible publisher page verifies the chapter's identity, but does not expose the claimed passage. This entry is retained as a motivated formal research programme, with that attribution marked unverified rather than asserted as a documented open conjecture.
\subsection*{References for the source of the problem}
Catherine Herfeld, \emph{Relating Neurons, Norms, and Networks to Rational Choice}, Chapter~5 of \emph{Conversations on Rational Choice}, Cambridge University Press, print year 2026, pp.~270--359; published online 27 December 2025. The chapter includes a conversation with Colin Camerer. \url{https://doi.org/10.1017/9781316344392.005}. Publisher record: \url{https://www.cambridge.org/core/books/abs/conversations-on-rational-choice/relating-neurons-norms-and-networks-to-rational-choice/48AC3B28E6FB44D73B84EA1207775605}. A precise habit-passage locator could not be independently verified.

\subsection*{Common conventions: Source mathematical conventions}

\textbf{Well-defined objects.}
All probability spaces and measurable variables are specified or inherited
from a local statistical model. Expectations used as finite targets require
the stated integrability. A decision rule or estimator is measurable with
respect to its available information; randomized rules may use an auxiliary
independent random variable. Norms without a subscript are Euclidean in
finite-dimensional spaces. An optimizer is requested only when attainment
is assured; otherwise the question concerns the optimal value and
$\varepsilon$-optimal admissible rules for every $\varepsilon>0$.

\textbf{Identification.}
A structural model $m\in\mathcal M$ implies an observable law
$P_m^{\rm obs}$ and target $\theta(m)$. The target is point identified on
$\mathcal M$ if
\[
 P_{m_1}^{\rm obs}=P_{m_2}^{\rm obs}
 \quad\Longrightarrow\quad\theta(m_1)=\theta(m_2)
 \qquad(m_1,m_2\in\mathcal M).
\]
When this fails, the sharp identified set is
\[
 \Theta(P;\mathcal M)=\{\theta(m):m\in\mathcal M,
                                  P_m^{\rm obs}=P\}.
\]
Specifying an intervention or estimand does not establish identification.
Positivity, exclusion, causal sufficiency, measurement restrictions, and
transport assumptions are substantive local inputs, never automatic
consequences of notation.

\textbf{Minimax risk and nontrivial inference.}
For a statistical experiment with data law $P^{(n)}$, target $\theta(P)$,
and nonnegative loss $\ell$, define
\[
 \mathcal R_n(\mathcal P,\ell)
   =\inf_{\widehat\theta_n}\sup_{P\in\mathcal P}
      \E_{P^{(n)}}\ell(\widehat\theta_n,\theta(P)).
\]
The infimum is over measurable estimators. A sharp rate is a sequence
$r_n$ for which $0<c\leq C<\infty$, independent of $n$, satisfies
$c r_n\leq\mathcal R_n\leq C r_n$ for all sufficiently large $n$.
Squared-error parametric risk is of order $n^{-1}$; the corresponding
root-mean-squared error is of order $n^{-1/2}$. These different conventions
are distinguished locally. A uniform asymptotic confidence region $C_n$
must satisfy
\[
 \liminf_{n\to\infty}\inf_{P\in\mathcal P}
     \Pr_{P^{(n)}}\{\theta(P)\in C_n\}\geq 1-\alpha,
 \qquad 0<\alpha<1,
\]
together with an informative diameter or expected-width requirement.
Coverage alone permits the vacuous whole parameter space. Testing questions
similarly impose both size control and power against a declared separated
alternative class.

\textbf{Binary causal baseline.}
When an entry explicitly invokes this baseline, one observes independent
copies of $O=(X,W,Y)$, with $W\in\{0,1\}$ and potential outcomes $Y(0),Y(1)$.
Assume consistency $Y=Y(W)$, no interference, conditional exchangeability
$(Y(0),Y(1))\perp W\mid X$, and overlap
$\epsilon\leq e(X)\leq1-\epsilon$ almost surely for a fixed
$0<\epsilon<1/2$, where
\[
 e(x)=\Pr(W=1\mid X=x),\qquad
 m_w(x)=\E[Y\mid W=w,X=x],\qquad
 \tau(x)=m_1(x)-m_0(x).
\]
These assumptions identify conditional mean effects almost everywhere
under the covariate law. Pointwise or uniform effects require appropriate
regularity and a specified support. Bounded outcomes, densities, and
smoothness are additional assumptions only where stated.

\textbf{Function classes.}
On $[0,1]^d$, $\mathcal H_d^s(L)$ is the isotropic Hölder ball of radius
$L>0$: put $k=\lceil s\rceil-1$ and $\gamma=s-k\in(0,1]$, bound derivatives
through order $k$, and require the order-$k$ derivatives to be
$\gamma$-Hölder with constant at most $L$. Thus integer orders use a
Lipschitz final derivative convention. The class
$\operatorname{BV}_d(L)$ consists of functions $f\in L^1((0,1)^d)$ whose
distributional derivative is a finite vector measure and for which
$\|f\|_{L^1}+|Df|((0,1)^d)\leq L$.
An $L^\infty$ bound or a particular pointwise version is imposed separately
when needed. Sparse and neural-network classes require a declared
dictionary or architecture and coefficient bounds.

An anisotropic H\"older class with declared block smoothness indices means
coordinatewise H\"older regularity in each block, uniformly in the other
blocks, with all corresponding derivative and H\"older bounds at most
the stated radius. Derivative orders and the integer-order convention in
each block are those just given. Mixed-derivative bounds are additional
restrictions only when explicitly stated.

\textbf{Decision and computational questions.}
Policy regret compares admissible feasible policies under the same law and
constraints. In computational entries, finite data and thresholds have
rational encodings; running time is measured in their total bit length.
Real-valued equilibrium search uses the explicitly declared algebraic or
FIXP encoding. An approximation ratio for maximization is written
$\operatorname{OPT}/\operatorname{ALG}$ and for minimization as
$\operatorname{ALG}/\operatorname{OPT}$, with zero optima and infeasible
instances handled locally. Historical approximation bounds are not
automatically current bounds.
% END ECONOMICS STATEMENT
\end{document}
